% ============================================================
% Content: More than Two States --- How a Qubit Encodes
% Superposition, registers, entanglement and qudits
% in the FFGFT framework
% Revised March 2026
% Builds on Doc.~173 (resonator foundations)
%           and Doc.~174 (spin manipulation, spin readout)
% Shor's algorithm: see Doc.~176
% ============================================================

\providecommand{\ket}[1]{\left|#1\right\rangle}
\providecommand{\bra}[1]{\left\langle#1\right|}
\providecommand{\braket}[2]{\left\langle#1\middle|#2\right\rangle}

% ============================================================
\section{The Real Question: What Lies Between \texorpdfstring{$\ket{0}$}{|0>}
  and \texorpdfstring{$\ket{1}$}{|1>}?}
\label{sec:between}
% ============================================================

Spin-up and spin-down are two points.
Two points are unambiguously labelled --- one bit of information.

But a \textbf{qubit} is not a bit.
The real power of the qubit lies not in
the fact that it can be $\ket{0}$ \emph{or} $\ket{1}$ ---
a classical coin can do that too.
It lies in the fact that it can be $\ket{0}$ \emph{and} $\ket{1}$
\emph{simultaneously} in a superposition ---
and in a continuous ratio:
\begin{equation}
  \ket{\psi} \;=\; \cos\!\tfrac{\theta}{2}\,\ket{0}
  \;+\; e^{i\phi}\sin\!\tfrac{\theta}{2}\,\ket{1}.
\end{equation}
Two parameters: the mixing angle $\theta \in [0,\pi]$
and the relative phase $\phi \in [0, 2\pi)$.
Together they parameterise a \textbf{sphere surface} ---
the \textbf{Bloch sphere} (geometric representation of
all possible single-qubit states as points
on a sphere of unit radius).

On this sphere there are \emph{infinitely many points} ---
not only the two poles $\ket{0}$ (north pole) and
$\ket{1}$ (south pole), but the entire equator,
all circles of latitude, every point in between.
Every point is a valid, physically
distinguishable quantum state.

\begin{keypoint}[A qubit does not encode 2 states --- it encodes \emph{infinitely many}]
The two classical poles $\ket{0}$ and $\ket{1}$
are special cases of a continuous family.
What distinguishes a qubit from a classical bit
is not the number of discrete states ---
it is the \textbf{continuous geometry} of the state space.

The readout interacts with the system --- every measurement
couples to the torsion configuration and drives it
into the nearest stable extreme state ($\ket{0}$ or $\ket{1}$).
The result is deterministic --- it depends on the exact
initial state which we do not fully know before measurement.
The apparent ``probability'' of standard QM describes
this epistemic ignorance, not fundamental randomness.

\textbf{The states themselves are discrete.} The continuity
lies in the phase space \emph{before} measurement ---
the readout always delivers one of the two stable poles.
\end{keypoint}

% ============================================================
\section{The Cylindrical Phase Space --- FFGFT Geometry of the Qubit}
\label{sec:cylinder}
% ============================================================

The Bloch sphere is the standard representation.
In FFGFT (Doc.~034) there is a
physically more direct alternative:
the \textbf{cylindrical phase space}.

A qubit state is described by
three geometric coordinates:
\begin{equation}
  (z,\; r,\; \theta),
\end{equation}
where:
\begin{itemize}
  \item $z \in [-1, +1]$: the ``height'' along the
    cylinder axis --- describes the proportion
    of the torsion configuration in the direction of $\ket{0}$ vs.\ $\ket{1}$.
    ($z = +1$: pure torsion type~$\ket{0}$; $z = -1$: pure torsion type~$\ket{1}$)
  \item $r \in [0, 1]$: the radius in the cross-section ---
    indicates the strength of the intermediate torsion,
    i.e.\ how far the state is from both poles.
    ($r = 0$: stable pole; $r = 1$: maximum intermediate torsion)
  \item $\theta \in [0, 2\pi)$: the azimuthal angle ---
    the \textbf{relative phase} between the two basis components.
\end{itemize}

The relation to the Bloch sphere is direct:
$z = \cos\theta_\text{Bloch}$ and $r = |\sin\theta_\text{Bloch}|$.
But the cylinder is not a geometric decoration ---
it makes gate operations into \emph{geometric
transformations} in real space:

\begin{center}
\resizebox{\textwidth}{!}{%
\begin{tabular}{p{2.5cm}p{5cm}p{5cm}}
\toprule
\textbf{Gate} & \textbf{Hilbert space (matrix)} &
\textbf{FFGFT cylinder (transformation)} \\
\midrule
\textbf{Hadamard H} &
$\frac{1}{\sqrt{2}}\begin{pmatrix}1&1\\1&-1\end{pmatrix}$ &
Swap $z \leftrightarrow r$, rotate phase by $\pi/2$ ---
moves pole to equator \\[8pt]
\textbf{Phase Z} &
$\begin{pmatrix}1&0\\0&-1\end{pmatrix}$ &
Add $\pi$ to phase $\theta$ ---
pure rotation about cylinder axis \\[8pt]
\textbf{Bit-flip X} &
$\begin{pmatrix}0&1\\1&0\end{pmatrix}$ &
2D rotation of $(z, r)$ by angle
$\alpha = \pi \cdot \Kfrak$ ---
fractal damping built in \\
\bottomrule
\end{tabular}
}
\end{center}

\begin{foundation}[FFGFT: gates as torsion rotations]
In the cylindrical FFGFT formalism every gate operation
is a geometric transformation of the torsion vector.
The Hadamard gate does not rotate a matrix ---
it tilts the torsion configuration from the
longitudinal axis to the transverse plane.
The Z gate rotates the phase $\theta$ ---
it rotates the torsion winding about its own axis.

The factor $\Kfrak$ in the X gate is not
a correction factor added afterwards ---
it follows directly from the fractal spacetime geometry
($\Dfrak = 2.999867$) and means:
every qubit rotation is slightly smaller than $\pi$,
because spacetime itself minimally deviates from flat
3D geometry.
This is a testable prediction:
all $\pi$ gates should systematically deviate by
$\Delta\alpha \approx 0.013\,\%$ from the ideal value.
\end{foundation}

% ============================================================
\section{Quantum Registers --- Exponential Growth through Combination}
\label{sec:register}
% ============================================================

A single qubit has infinitely many intermediate states
in the torsion phase space, but always delivers one of
the two stable poles upon readout --- the readout itself
is the interaction that drives the state to the next
stable extreme. The real power comes through \textbf{combination}.

\subsection{Product States: Additive Classicality}

Two independent qubits have together
$2 \times 2 = 4$ basis states:
$\ket{00}$, $\ket{01}$, $\ket{10}$, $\ket{11}$.

The general state \emph{without entanglement}
(product state) is:
\begin{equation}
  \ket{\psi_A} \otimes \ket{\psi_B}
  \;=\; (\alpha_A\ket{0} + \beta_A\ket{1})
         \otimes
         (\alpha_B\ket{0} + \beta_B\ket{1}).
\end{equation}
This is only two separate Bloch spheres ---
no gain over two independent qubits.
The state space is \emph{additive}: $2 + 2 = 4$ parameters.

\subsection{Entangled States: Multiplicative Growth}
\label{sec:entanglement}

Once the two qubits are \emph{entangled},
the joint state can no longer be written as a product
of two individual states.
The general two-qubit state requires
$4$ complex amplitudes (minus normalisation and
global phase: $2 \times 4 - 2 = 6$ real parameters):
\begin{equation}
  \ket{\Psi} \;=\; \alpha\ket{00} + \beta\ket{01}
              + \gamma\ket{10} + \delta\ket{11}.
\end{equation}
This is \textbf{more} than two independent qubits ---
entangled qubits carry joint information
that is contained in neither individual qubit.

For $n$ qubits the Hilbert space grows exponentially:
\begin{equation}
  \boxed{d(n) \;=\; 2^n}
\end{equation}
Real parameters in the general state: $2 \cdot 2^n - 2$.

\begin{center}
\resizebox{\textwidth}{!}{%
\begin{tabular}{rrrr}
\toprule
$n$ qubits & Basis states & Real parameters & Classical equivalent \\
\midrule
1  & 2            & 2         & 1 bit \\
2  & 4            & 6         & 2 bits \\
3  & 8            & 14        & 3 bits \\
10 & 1,024        & 2,046     & 10 bits \\
50 & $10^{15}$    & $2\times10^{15}$ & 50 bits \\
300 & $10^{90}$   & $2\times10^{90}$ & $\approx$ number of atoms in the universe \\
\bottomrule
\end{tabular}
}
\end{center}

\textbf{This is the actual quantum advantage:}
Not a single qubit is powerful ---
but the collective growth of an
\emph{entangled} register.

\subsection{Bell States --- Maximum Entanglement in Two Qubits}
\label{sec:bell}

The four \textbf{Bell states} are the
maximally entangled two-qubit states ---
each a complete orthonormal basis
of the entangled subspace:
\begin{align}
  \ket{\Phi^+} &= \tfrac{1}{\sqrt{2}}(\ket{00} + \ket{11}), \\
  \ket{\Phi^-} &= \tfrac{1}{\sqrt{2}}(\ket{00} - \ket{11}), \\
  \ket{\Psi^+} &= \tfrac{1}{\sqrt{2}}(\ket{01} + \ket{10}), \\
  \ket{\Psi^-} &= \tfrac{1}{\sqrt{2}}(\ket{01} - \ket{10}).
\end{align}
Each of these states has the property:
if qubit~A is measured, the result of qubit~B
is immediately determined --- regardless of how far away.

In FFGFT (Doc.~023, 035):
\begin{equation}
  E^{\text{FFGFT}}(a, b) \;=\;
  -\cos(a - b) \cdot
  \left(1 - \xi\,\frac{(\Delta\theta)^2/\pi^2}{\Dfrak}\right),
\end{equation}
the $\xi$ damping reduces the CHSH violation from
$2\sqrt{2} \approx 2.8284$ to
$\approx 2.8275$ --- a difference of
$\sim 10^{-4}$, which explains non-locality locally
without eliminating it.

\begin{foundation}[FFGFT: Bell states as joint torsion nodes]
In the FFGFT picture a Bell state is not an
action at a distance between two qubits ---
it is a \textbf{topologically connected
torsion configuration} across two resonators.
The two resonators share a common
torsion node whose resolution
simultaneously determines both ends
at every local measurement.

This is not a signal propagating faster than light ---
it is a common geometric state established at preparation.
The local readout interaction at one end
deterministically drives the local torsion into a
stable pole --- and thereby necessarily also
determines the state of the other end, because both share the same
torsion node.
The ``global consequence'' is not action at a distance,
but the geometric coherence of the common node.

The $\xi$ damping (FFGFT correction to Bell correlations)
is in this picture the finite extension of the
torsion node: the node is not point-like,
but has a characteristic length
$\sim L_\xi = \xi \cdot \ell_P \cdot (1/\xi)^N$
at the $\xi$ scale level $N \approx 8$.
\end{foundation}

\begin{keyresult}[IBM Quantum hardware validation: Bell states, June 2025]
The FFGFT predictions for Bell states were experimentally
verified on genuine quantum hardware:

\textbf{Hardware:} IBM Brisbane \& IBM Sherbrooke (127 qubits each,
Heavy-Hex lattice), access via IBM Quantum Platform / Qiskit Runtime.

\textbf{Result:}
\begin{center}
{\footnotesize
\begin{tabularx}{\textwidth}{lX}
\toprule
Quantity & Value \\
\midrule
Bell state fidelity        & $97.17\,\%$ \\
FFGFT algorithm vs.\ QM   & Agreement within measurement precision \\
Repetition variance        & $0.0001$--$0.001$ (excellent) \\
Deviation $P(00)$, $P(11)$ & $1$--$5\,\%$ (hardware noise) \\
\bottomrule
\end{tabularx}
} % footnotesize
\end{center}

\textbf{Significance in FFGFT context:}
The FFGFT amplitudes $E_i(x,t)$ with $\xi$ coupling
($\xi = 1.038 \times 10^{-5}$, Higgs-derived) deliver
algorithmically equivalent predictions to standard QM ---
with improved repeatability (lower variance).
The $\xi$ deviation from the ideal Bell correlation
($\Delta\text{CHSH} \approx 10^{-5}$) lies below
the current hardware noise threshold and could therefore
not yet be directly isolated and measured.
This is a precise experimental task for future
high-precision Bell tests.
\end{keyresult}

% ============================================================
\section{Three Ways to More States --- a Comparison}
\label{sec:three-ways}
% ============================================================

\begin{center}
\resizebox{\textwidth}{!}{%
\begin{tabular}{p{2.8cm}p{3cm}p{3.5cm}p{3.5cm}}
\toprule
\textbf{Strategy} & \textbf{More states through} &
\textbf{Growth} & \textbf{FFGFT meaning} \\
\midrule
\textbf{Superposition}\newline (1 qubit) &
Continuous superposition of $\ket{0}$ and $\ket{1}$ &
$\infty$, but 1 bit at measurement &
Intermediate torsion on the Bloch sphere \\[8pt]
\textbf{Register}\newline ($n$ qubits) &
Entanglement of multiple qubits &
$2^n$ (exponential) &
$2^n$-dimensional torsion product space; entangled = topologically connected nodes \\[8pt]
\textbf{Qudit}\newline ($d > 2$, 1 object) &
Higher energy levels, orbital states, exciton number &
$d$ (linear) per object, $d^n$ for $n$ qudits &
$d$ stable torsion configurations in one resonator \\
\bottomrule
\end{tabular}
}
\end{center}

\begin{keypoint}[No strategy dominates in all areas]
Register entanglement scales most strongly ($2^n$),
but requires $n$ physical objects and error-corrected
entanglement operations between them.

Qudit coding gains information \emph{per object}
($\log_2 d > 1$), but the exponential
growth is slower.

Superposition is always present --- it is the basis
for quantum interference, the actual driving force
of quantum algorithms.
\end{keypoint}

% ============================================================
\section{Quantum Interference --- Why Superposition is Useful}
\label{sec:interference}
% ============================================================

Superposition alone would be a mere overlay ---
the readout would always deliver one of the two stable poles,
deterministically dependent on the exact phase state.

What makes superposition into a \emph{resource} is
\textbf{interference}: amplitudes can
constructively (amplify) or destructively (cancel)
overlap --- depending on phase $\phi$.

This is the principle behind all useful
quantum algorithms:
\begin{enumerate}
  \item Prepare the register in a
    uniform intermediate torsion over all inputs
    (with the Hadamard gate).
  \item Apply a function that increases the phases
    of the correct answers and decreases those
    of the wrong ones.
  \item Let interference drive the wrong answers
    into unstable torsion configurations
    and the correct one into a stable one.
  \item The readout drives the system deterministically
    into the nearest stable pole ---
    which after interference corresponds
    with overwhelming frequency to the sought answer.
\end{enumerate}

Without interference: Monte Carlo (classical guessing).
With interference: Grover ($\sqrt{N}$ instead of $N$ steps),
Shor (factorisation in polynomial time).

\begin{foundation}[FFGFT: interference as torsion phase superposition]
In FFGFT the phase $\phi$ is not an abstract
complex number --- it is the rotation position
of the intermediate torsion configuration (Doc.~174).

Constructive interference: two torsion phases
coincide --- they amplify each other into a
geometrically more stable configuration.
Destructive interference: two opposite
torsion phases cancel each other ---
the resulting configuration has no
stable geometric anchor and is not addressed at
the readout interaction.

A quantum algorithm in FFGFT language:
a sequence of torsion rotations that aligns the
phase configuration of the register so that
the sought answer ends in a stable
torsion configuration --- all others
in geometrically unstable ones.
The readout is then not a random selection,
but a deterministic reaction of the
resonator to the readout interaction.
\end{foundation}

% ============================================================
\section{QND Measurement --- Reading Without Destroying}
\label{sec:qnd}
% ============================================================

\subsection{What QND Means}

A QND measurement reads an observable $\hat{A}$
without changing the corresponding eigenstate.
The formal condition:
\begin{equation}
  [\hat{A},\; \hat{H}_\text{measurement}] \;=\; 0.
\end{equation}
The measurement operator $\hat{H}_\text{measurement}$ commutes
with the measured observable --- it does not change the
eigenvalue, because it respects the same eigenspace.

Practically this means: the same state can be
\textbf{repeatedly measured} and delivers
the same result each time --- because the measurement
does not touch the state itself.

The most important example from Doc.~174:
\textbf{Faraday/Kerr rotation}.
A linearly polarised laser beam passes the
quantum dot without being absorbed ---
it rotates its plane of polarisation by $\theta_F \propto S_z$.
No photon is absorbed. No energy transfer
into the spin mode. The spin state before and after
measurement is identical.

\subsection{QND as Empirical Proof of Determinism}
\label{sec:qnd-proof}

The QND measurement is from the FFGFT perspective far more
than a technically useful property ---
it is an \textbf{empirical proof}
for the reality of stable torsion configurations.

The argument is direct:

\begin{enumerate}
  \item Faraday rotation measures $S_z$ without
    changing the spin state.
  \item Repeated measurements always yield the same result.
  \item Therefore the state had this value
    \emph{before every measurement} --- definitely,
    not as an abstract superposition that only
    ``becomes reality'' through measurement.
  \item The state was real before we looked.
\end{enumerate}

\begin{foundation}[FFGFT: QND measurement and stable torsion configuration]
In FFGFT, Faraday rotation is state-preserving
because a stable torsion configuration
(spin-up = topologically fixed winding in one direction)
cannot be driven out of its state through a dispersive,
non-resonant coupling.

The photon ``sees'' the torsion --- it rotates its
polarisation accordingly --- but it transfers no
energy into the torsion mode, because it is not resonant.
The geometric state is invariant under this coupling.

\textbf{The philosophical consequence is clear:}
If a state can be repeatedly measured without disturbance,
then it was not ``created'' by the measurement ---
it was there before.
The standard QM interpretation
(``the state exists only through measurement'')
is not compatible with QND measurements.

In FFGFT: the torsion configuration is a
geometric reality --- stable, definite, pre-existent.
The measurement reads it off. It does not create it.
\end{foundation}

\subsection{Potential Landscape: Stable Minima, Transitions and Readout Strength}

\begin{foundation}[FFGFT: torsion potential, transitions and readout strength]
The complete picture of qubit states is a
\textbf{potential landscape with two stable minima}
and a saddle between them.

\textbf{The readout strength decides what happens:}

\begin{center}
{\footnotesize\resizebox{\textwidth}{!}{%
\begin{tabular}{p{3cm}p{4.5cm}p{5cm}}
\toprule
\textbf{Initial state} & \textbf{Weak coupling} & \textbf{Strong coupling} \\
\midrule
Stable minimum\newline ($\ket{\uparrow}$ or $\ket{\downarrow}$) &
Coupling does not lift state from minimum ---
\textbf{QND}: state remains, result reproducible &
Coupling can lift state from minimum ---
invasive, state changed \\[8pt]
Slope / saddle\newline (superposition) &
Even weak coupling is sufficient to drive the state ---
it ``rolls'' into the next minimum.
\textbf{Which one}: determined by exact phase $\phi$
(epistemically unknown, deterministic) &
Same, just faster ---
state jumps immediately into next minimum \\
\bottomrule
\end{tabular}
}} % resizebox
\end{center}

\textbf{The decisive insight:}
There is no principled boundary between
``QND-possible'' and ``QND-impossible'' ---
there are only \textbf{more stable and less stable states}
in a continuous potential landscape.

A state in the deep minimum is stable against
the readout coupling --- the minimum must be overcome.
A state on the slope is unstable ---
any coupling, however weak, is sufficient to drive it
into the next minimum.

And the ``jump'' to the final position is not random ---
it is determined deterministically by the phase of the initial state
on the slope:
a state near $\ket{\uparrow}$ lands in $\ket{\uparrow}$,
a state near $\ket{\downarrow}$ lands in $\ket{\downarrow}$,
a state on the saddle ($\theta = \pi/2$) lands
depending on the fine phase $\phi$.

The apparent ``randomness'' of standard QM
is the epistemic ignorance of this phase ---
not an ontological randomness of the system.

The picture applies to all $d$ levels of a qudit:
not two minima, but $d$ minima ---
with $d-1$ saddle points between them.
\end{foundation}

\begin{keypoint}[QND and superposition --- no contradiction]
A stable eigenstate ($\ket{\uparrow}$ or $\ket{\downarrow}$):
sits in the deep minimum --- weak coupling reads
without driving. QND possible. Result reproducible.

An intermediate torsion ($\alpha\ket{\uparrow}+\beta\ket{\downarrow}$):
sits on the slope --- any coupling drives
into the next minimum. Deterministic, not random.
The phase $\phi$ determines which minimum is reached.

The QND measurement \emph{proves} the reality of stable poles.
The disturbing measurement \emph{proves} the instability
of intermediate torsions --- and the determinism structure
of the potential landscape.
Together: completely deterministic,
geometrically founded picture.
\end{keypoint}

% ============================================================
\section{Two Levels of the \texorpdfstring{$\xi$}{xi}-Hierarchy:
  Spin Qubit and Photonic Resonator}
\label{sec:two-levels}
% ============================================================

\subsection{Spin Qubit: Minimal Resonator at Level \texorpdfstring{$N \approx 8$}{N≈8}}

A spin qubit is the absolute minimum of a
quantum resonator (Doc.~173):
a single electron, a single energy level,
two permitted orientations.

The energy scale comes from
\textbf{quantum confinement} at $\xi$ level $N \approx 8$
(nanometre scale).
The spacing between the two spin states
in a typical magnetic field ($B \approx 1$\,T)
lies at:
\begin{equation}
  \Delta E_Z \;=\; g^* \mu_B B \;\approx\; 0.1\text{--}1\,\text{meV}.
\end{equation}
The thermal energy at room temperature is:
\begin{equation}
  k_B T(300\,\text{K}) \;=\; 26\,\text{meV}.
\end{equation}
So $\Delta E_Z \ll k_BT$ --- thermal noise
constantly knocks the system out of the minimum.
Low temperatures ($T < 1$\,K, $k_BT < 0.1$\,meV)
are not optional, but mandatory ---
not to maintain superposition,
but simply to \emph{stay} in one of the two minima.

\subsection{Photonic Resonator: Collective Geometry at Level
  \texorpdfstring{$N \approx 9$--10}{N≈9-10}}

A waveguide of lithium niobate (LiNbO$_3$),
a Mach-Zehnder interferometer, a ring resonator ---
these are resonators made of $10^{18}$ to $10^{23}$ atoms
in perfect lattice order.

The decisive difference from the spin qubit:
\textbf{the resonator is not a single particle ---
it is the macroscopic geometry of the entire crystal.}

The energy scale of an optical photon:
\begin{equation}
  E_\text{photon} \;\approx\; 1\text{--}2\,\text{eV}
  \quad (780\text{--}1550\,\text{nm}),
\end{equation}
so:
\begin{equation}
  \frac{E_\text{photon}}{k_BT(300\,\text{K})}
  \;\approx\; \frac{1\,\text{eV}}{26\,\text{meV}}
  \;\approx\; 40.
\end{equation}
Thermal noise cannot spontaneously create optical
photons --- the system sits far below
the thermal threshold for optical excitations.
Condition~2 from Doc.~173 ($\Delta E \gg k_BT$)
is satisfied by a factor of~40.

\begin{important}[Why photonics works at room temperature]
Not because the photonic system is more robustly built ---
but because it operates at a higher $\xi$ level
of the energy hierarchy.

The \textbf{geometric robustness} (Condition~1)
is extraordinarily well satisfied by the macroscopic
single crystal with $\sim 10^{23}$ lattice sites in perfect order ---
much more robust than a single quantum dot.
Individual phonons (thermal lattice vibrations)
barely disturb the collective mode.

The \textbf{thermal robustness} (Condition~2)
is automatically given by the eV energy scale ---
thermal photons in the infrared have $\sim 25$\,meV,
optical photons have $\sim 1$\,eV:
four orders of magnitude safety margin.
\end{important}

\subsection{Comparison of the Two Levels in the \texorpdfstring{$\xi$}{xi}-Hierarchy}
\label{sec:xi-comparison}

\begin{center}
\resizebox{\textwidth}{!}{%
\begin{tabular}{p{3.5cm}p{5cm}p{5cm}}
\toprule
& \textbf{Spin qubit} & \textbf{Photonic resonator} \\
\midrule
\textbf{$\xi$ level} &
$N \approx 8$ (nm scale) &
$N \approx 9$--10 ($\mu$m--mm scale) \\[4pt]
\textbf{Resonator} &
1 electron in nm cage &
$10^{18}$--$10^{23}$ atoms as collective crystal \\[4pt]
\textbf{Energy scale} &
$\Delta E \approx 0.1$--1\,meV &
$E_\text{photon} \approx 1$--2\,eV \\[4pt]
\textbf{$\Delta E / k_BT$(300\,K)} &
$\approx 0.004$--0.04 (too small!) &
$\approx 40$ (robust) \\[4pt]
\textbf{Cooling needed?} &
Yes, $T < 1$\,K mandatory &
No, room temperature \\[4pt]
\textbf{Analogue states} &
Fragile: $T_2 \sim$ ns at 300\,K &
Stable: phase thermally unreachable \\[4pt]
\textbf{Discreteness from} &
Quantum confinement (1 particle) &
Macroscopic mode geometry \\[4pt]
\textbf{Typical application} &
Digital qubit, quantum gate &
Analogue signal processing, photonic computing, 6G \\
\bottomrule
\end{tabular}
}
\end{center}

\begin{foundation}[FFGFT: two levels of the $\xi$ hierarchy ---
  same physics, different scales]
Spin qubit and photonic resonator are
in FFGFT not fundamentally different systems ---
they are resonators at different levels
of the $\xi$ scaling ladder.

At level $N \approx 8$: the confinement energy minimum
lies at meV --- not thermally robust at 300\,K.
At level $N \approx 9$--10: the collective mode minimum
lies at eV --- completely thermally robust.

The $\xi$ parameter generates a universal
energy hierarchy: each level carries the energy
by a factor of $1/\xi \approx 7500$ higher than the previous.
From meV (level~8) to eV (level~9--10):
the factor is $\sim 10^3$--$10^4$ ---
consistent with $\xi$ scaling.

This explains not only photonics at room temperature ---
it explains why biological systems use photonic
mechanisms (photosynthesis, bird navigation via light)
and not spin qubit systems for quantum effects:
nature works at the thermally robust level.
\end{foundation}

% ============================================================
\section{\texorpdfstring{$\xi$}{xi}-Level Separation as Stability Principle}
\label{sec:xi-separation}
% ============================================================

\subsection{The Coupling Formula Between Levels}

Between different $\xi$ levels, the coupling strength decreases
exponentially with the level difference $\Delta N$:
\begin{equation}
  \Gamma_{\Delta N} \;\propto\; \xi^{\Delta N}.
\end{equation}
For $\Delta N = 1$: $\xi^1 \approx 1.333 \times 10^{-4}$ (weak but non-zero).
For $\Delta N = 2$: $\xi^2 \approx 1.78 \times 10^{-8}$ (extremely weak).
For $\Delta N = 4$: $\xi^4 \approx 3.2 \times 10^{-16}$ (practically non-existent).

\subsection{Fractal Self-Similarity \emph{Because of} the Damping}

\textbf{The fractal hierarchy of nature is not stable
\emph{despite} $\xi$ damping --- it is stable
\emph{because of} it.}

Were $\xi$ not small, all levels would
strongly couple to each other.
Every change at one level would immediately
restructure all other levels.
There would be no stable atoms, no stable molecules,
no stable cells --- only a single
chaotic coupled system of all scales simultaneously.

Because $\xi \ll 1$, the levels can
exist \textbf{quasi-independently}:
each level has its own stable internal dynamics,
barely disturbed by the other levels.

\begin{foundation}[FFGFT: $\xi$ damping as origin of
  the natural hierarchy]
The existence of separate, stable structural levels
in nature --- quarks, nucleons, atoms,
molecules, cells, organisms, planets, galaxies ---
is not a random property of matter.

It follows directly from $\xi$ level separation:
each level is practically invisible to the next-but-one.
The coupling falls exponentially with level distance.
This is the mechanism that makes the fractal
self-similarity of nature \emph{stable}.
\end{foundation}

% ============================================================
\section{Qudit Register --- the Connection to the \texorpdfstring{$\xi$}{xi}-Hierarchy}
\label{sec:qudit-register}
% ============================================================

From Doc.~174: a quantum dot of dimension $d$
encodes $\log_2 d$ bits.
What happens when one entangles $n$ such qudits?

The total state space has dimension $d^n$.
For $d = 4$ (spin-orbital ququart) and $n$ qudits:
\begin{equation}
  d^n \;=\; 4^n \;=\; (2^2)^n \;=\; 2^{2n}.
\end{equation}
This is equivalent to $2n$ qubits in a register
of size $n$ --- double the information density per object.

\begin{keyresult}[Qudit density at $\xi$ level $N \approx 8$]
A quantum dot with $R = 3$\,nm at 4\,K
has at least 4 thermally stable states
(spin-orbital ququart, $d = 4$).
A register of $n = 10$ such qudits
has dimension $4^{10} = 1,048,576$ ---
the same as 20 classical qubits, but
only 10 physical objects at $\xi$ level $N \approx 8$.
\end{keyresult}

% ============================================================
\section{The FFGFT Topology Compiler --- Geometry of Entanglement}
\label{sec:compiler}
% ============================================================

A practical problem with entangled registers:
entanglement gates between non-adjacent qubits
are expensive --- they require many SWAP operations
(which accumulate errors and cost time).

FFGFT (Doc.~034) offers a geometric
solution: the \textbf{FFGFT topology compiler}.

Standard compilers minimise the number of SWAPs.
The FFGFT compiler instead minimises the cumulative
\textbf{fractal path length} of all entanglement operations.
Since the fractal damping term $\exp(-\xi \cdot \ell / \Dfrak)$
is distance-dependent (larger physical distance $\ell$
means stronger $\xi$ damping of coherence),
critically interacting qubits must be
\emph{physically placed closer} to each other.

This leads to a new chip architecture:
not laid out according to logical circuit optimisation,
but according to \textbf{geometric coherence maximisation}.

\begin{significance}[Connection to the adaptive hybrid architecture]
The adaptive hybrid architecture (Doc.~173 conclusion)
selects resonators according to task:
quantum dots for local synaptic operations,
Ising machines for global optimisation.

The FFGFT topology compiler is the linking element:
it decides which qubits/qudits must physically
sit close together (strong entanglement)
and which may remain further apart
(weakly, tolerantly damped entanglement).

The learning AI on this hardware does not only learn
\emph{what} is to be computed ---
it also learns \emph{which geometric arrangement}
executes the computation most coherently.
Not through programmed topology specifications,
but through geometrically informed learning
in the FFGFT compiler.
\end{significance}

% ============================================================
\section{Comparison: Classical Bit, Qubit, Qudit, Register}
\label{sec:comparison}
% ============================================================

\begin{center}
\resizebox{\textwidth}{!}{%
\begin{tabular}{p{2.5cm}p{2.5cm}p{2.5cm}p{2.5cm}p{2.5cm}}
\toprule
& \textbf{Class.\ bit} & \textbf{Qubit} &
\textbf{Qudit ($d$)} & \textbf{$n$ qubits} \\
\midrule
\textbf{State space} &
2 points &
Bloch sphere ($\infty$) &
$\mathbb{CP}^{d-1}$ ($\infty$) &
$\mathbb{CP}^{2^n-1}$ ($\infty$) \\[4pt]
\textbf{Readout result} &
1 bit (deterministic) &
1 bit (deterministic,\newline depending on phase state) &
$\log_2 d$ bits (deterministic) &
$n$ bits \\[4pt]
\textbf{Parallel processing} &
none &
Interference over 2 paths &
Interference over $d$ paths &
Interference over $2^n$ paths \\[4pt]
\textbf{Resource for advantage} &
--- &
Phase $\phi$ &
$d-1$ phases &
Entanglement + phases \\[4pt]
\textbf{FFGFT picture} &
stable/unstable torsion &
Torsion on Bloch sphere &
$d$ stable torsion levels &
Torsion network, topologically connected \\
\bottomrule
\end{tabular}
}
\end{center}

% ============================================================
\section{Open Connections --- Next Steps}
\label{sec:open-175}
% ============================================================

\begin{significance}[Open questions for further chapters]
\textbf{1. Fractal gate corrections measurable?}\\
The FFGFT compiler says: all $\pi$ pulses deviate
by $\Delta\alpha \approx \Kfrak - 1 \approx 0.013\,\%$
from the ideal $\pi$.
Current superconducting qubit systems achieve
gate fidelity $>99.9\,\%$.
That is just at the boundary to see this
systematic deviation.
Is there a systematic bias in the $\pi$ pulse length
in current calibration data?\\[6pt]

\textbf{2. Bell states and $\xi$ scale level:}\\
The $\xi$ damping in Bell correlations
($\Delta\text{CHSH} \approx 10^{-4}$) depends on $N \approx 8$.
For Bell tests with \emph{macroscopic} objects
(atoms, molecules) the scale level $N$ would increase ---
and with it the damping.
Is there an $N$-dependent suppression
of Bell violations for different system sizes?\\[6pt]

\textbf{3. Qudit entanglement and torus topology:}\\
Two entangled qudits ($d = 4$) correspond
in FFGFT to two topologically connected
spin-orbital torsion configurations.
What is the geometric description
of this connection in the toroidal field?
Is there a conserved quantity (topological invariant)
that quantifies the entanglement depth?\\[6pt]

\textbf{4. Quantum error correction and $\xi$ granularity:}\\
The surface code needs $\sim 1000$ physical qubits
per logical qubit.
In a qudit-based register with $d = 4$
one needs only half as many physical objects
for the same logical information.
How does the error correction threshold change
in the FFGFT framework when the natural noise floor
$\xi \approx 1.333 \times 10^{-4}$ applies as
the lower error bound?
\end{significance}

% ============================================================
\section{Summary}
\label{sec:summary-175}
% ============================================================

\begin{conclusionBox}[Conclusion: More than Two States]
A single qubit has infinitely many intermediate states
in the torsion phase space --- but the states themselves are
\textbf{discrete}: the readout always delivers one of the
two stable poles, deterministically determined by the
exact phase state at the moment of readout.
The apparent ``probability'' of standard QM
is epistemic ignorance of this state ---
not fundamental randomness of the system.

\textbf{The readout interacts with the system ---
but not always destructively.}
Faraday/Kerr rotation (Doc.~174) is a
QND measurement: it reads the state without changing it,
because the coupling commutes with the spin operator.
That the same state can be repeatedly measured
is the empirical proof that it \emph{existed before} measurement ---
definitively and really.
The result remains deterministic in every case.

The actual quantum advantage comes through three routes:

\textbf{1. Superposition:}
Continuous intermediate torsion with phase $\phi$
enables interference --- the geometric foundation
of all useful quantum algorithms.
Interference aligns the register so that
the sought answer ends in a stable torsion configuration;
the readout follows that deterministically.

\textbf{2. Register entanglement:}
$n$ qubits span a $2^n$-dimensional torsion space.
Entanglement means: the torsion nodes of the qubits
are topologically connected --- reading one
node necessarily also determines the others.
Bell correlations are locally explained by $\xi$ damping
without eliminating non-locality.

\textbf{3. Qudit coding:}
A resonator with $d > 2$ stable torsion levels
delivers $\log_2 d > 1$ bits per object at readout.
Spin-orbital ququarts ($d = 4$) at
$\xi$ level $N \approx 8$ double the
information density of a register.

In FFGFT all three strategies are expressions of the same
fundamental principle: stable torsion configurations
in geometrically precise resonators --- discrete from inside,
analoguely controllable from outside, deterministic at readout.
\end{conclusionBox}

% ============================================================

\begin{thebibliography}{99}

\bibitem{Kouwenhoven1997}
Kouwenhoven, L.\,P. et al.
\textit{Excitation spectra of circular, few-electron quantum dots.}
Science \textbf{278}, 1788--1792 (1997).

\bibitem{Loss1998}
Loss, D. \& DiVincenzo, D.\,P.
\textit{Quantum computation with quantum dots.}
Phys.\ Rev.\ A \textbf{57}, 120--126 (1998).

\bibitem{Veldhorst2014}
Veldhorst, M. et al.
\textit{An addressable quantum dot qubit with fault-tolerant control-fidelity.}
Nature Nanotechnol.\ \textbf{9}, 981--985 (2014).

\bibitem{Muhonen2014}
Muhonen, J.\,T. et al.
\textit{Storing quantum information for 30 seconds in a nanoelectronic device.}
Nature Nanotechnol.\ \textbf{9}, 986--991 (2014).

\bibitem{Atature2007}
Atatüre, M. et al.
\textit{Observation of Faraday rotation from a single confined spin.}
Nature Phys.\ \textbf{3}, 101--105 (2007).

\bibitem{Grangier1998}
Grangier, P., Levenson, J.\,A. \& Poizat, J.\,P.
\textit{Quantum non-demolition measurements in optics.}
Nature \textbf{396}, 537--542 (1998).

\bibitem{Wang2018}
Wang, C. et al.
\textit{Integrated lithium niobate electro-optic modulators operating at
CMOS-compatible voltages.}
Nature \textbf{562}, 101--104 (2018).

\bibitem{Reck1994}
Reck, M. et al.
\textit{Experimental realization of any discrete unitary operator.}
Phys.\ Rev.\ Lett.\ \textbf{73}, 58--61 (1994).

\bibitem{Shannon1948}
Shannon, C.\,E.
\textit{A mathematical theory of communication.}
Bell Syst.\ Tech.\ J.\ \textbf{27}, 379--423 (1948).

\bibitem{Knill2001}
Knill, E., Laflamme, R. \& Milburn, G.\,J.
\textit{A scheme for efficient quantum computation with linear optics.}
Nature \textbf{409}, 46--52 (2001).

\end{thebibliography}
